% !TeX root = ../main.tex

\xchapter{相关理论和知识背景}{Related Theories and Background Knowledge}

\xsection{基于 Wi-Fi 的 CSI 信号的人类活动识别原理}{Principles of Human Activity Recognition Based on Wi-Fi CSI Signals}

如第一章所述，基于 Wi-Fi 信道状态信息（Channel State Information, CSI）的人类活动识别通过分析人体活动对无线信号传播的影响实现动作推断。本节从信号层面阐述 CSI 的物理含义与数据表示，为第三、四章的模型输入提供技术背景。图~\ref{fig:har-csi-flow} 给出了基于 Wi-Fi CSI 信号的人类活动识别系统通用流程，其核心是建立人体行为与 CSI 信号变化之间的映射关系。

\begin{figure}[H]
  \centering
  \includegraphics[width=0.85\textwidth]{fig2-1-har-csi-flow.png}
  \caption{基于 Wi-Fi CSI 的人类活动识别系统通用流程}
  \label{fig:har-csi-flow}
\end{figure}

CSI 是描述无线信道传播特性的核心参量，刻画了信号在发射端与接收端之间沿各条路径传播时的幅度衰减与相位偏移。在室内环境中，无线信号经反射、折射、散射、衍射等多径效应到达接收端，各路径对最终接收信号产生不同的幅度与相位贡献，CSI 正是对这些变化的精细刻画\cite{ma2019wifi}。因此，当人体进入传播区域时，其移动、遮挡、反射等会改变多径结构，进而引起 CSI 的时变特征，为活动识别提供了物理依据。

在遵循 IEEE 802.11 标准的现代 Wi-Fi 系统中，物理层通过多输入多输出（Multiple-Input Multiple-Output, MIMO）与正交频分复用（Orthogonal Frequency-Division Multiplexing, OFDM）协同工作。MIMO 在收发两端配置多天线，利用多径效应提升吞吐量；OFDM 将宽带信号划分为若干窄带子载波，各子载波在频域正交、互不干扰。对于 MIMO-OFDM 信道，每个子载波上每对收发天线对应一个 CSI 复数响应。信道频响（Channel Frequency Response, CFR）可表示为
\begin{equation}
  H(f; t) = \sum_{n=1}^{N} a_n(t) \mathrm{e}^{-\mathrm{j}2\pi f \tau_n(t)}
  \tag{2-1}
\end{equation}
式中：$a_n(t)$ 为第 $n$ 条路径的幅度衰减因子；$\tau_n(t)$ 为传播时延；$f$ 为载波频率。CSI 的幅度 $|H|$ 与相位 $\angle H$ 受发射端、接收端及环境中物体位移与运动影响，从而将环境变化蕴含于信道特性之中。

对于具有 $M$ 根发射天线、$N$ 根接收天线、$K$ 个子载波的 MIMO-OFDM 系统，单次测量的 CSI 可表示为三维复数矩阵 $H \in \mathbb{C}^{N \times M \times K}$，分别对应空间维度（收发天线对）与频域维度（子载波）。在时间维度上对 CSI 进行连续采样，可得到四维 CSI 张量，在时、频、空三个域内刻画信道随人体活动产生的变化。CSI 数据在子载波维与数据包（或时间戳）维上组织，每个时间戳对应一个 $N_{\mathrm{T}} \times N_{\mathrm{R}} \times N_{\mathrm{sub}}$ 的矩阵，其中 $N_{\mathrm{T}}$、$N_{\mathrm{R}}$、$N_{\mathrm{sub}}$ 分别为发射天线数、接收天线数与子载波数。

在 IEEE 802.11n/ac 协议栈中，CSI 的获取通常依赖物理层前导结构。发射端在每帧前导内插入长训练域（Long Training Field, LTF）或高吞吐量长训练域（HT-LTF），接收端根据已知的发送序列与接收信号估计各子载波的信道响应，进而得到 CSI 矩阵。商用网卡（如 Intel 5300）通过修改驱动或固件可导出每根天线、每个子载波上的 CSI 至用户态，使得基于 CSI 的感知应用得以在通用硬件上实现，无需额外传感器。原始 CSI 数据常含有噪声、载波间干扰及硬件偏差，因此在特征提取前需进行预处理：常见步骤包括异常值剔除、线性趋势去除、滑动平均或滤波以抑制高频噪声，以及必要时对幅度与相位进行标准化或归一化，为后续特征提取与分类奠定基础。

相较接收信号强度指示（Received Signal Strength Indication, RSSI），CSI 具有更高的分辨能力。RSSI 仅提供单一标量功率值，而 CSI 在多个子载波与天线对上提供复数响应，能更细致地反映多径结构，因而更适合细粒度感知任务。在室内活动识别中，CSI 幅度的时域变化可区分不同人体、动作与手势，相位的空频域变化可用于定位与呼吸估计等\cite{ma2019wifi}。信道冲激响应 $h(\tau)$ 在频域的另一种表述形式为
\begin{equation}
  h(\tau) = \sum_{l=1}^{L} \alpha_l \mathrm{e}^{\mathrm{j}\varphi_l} \delta(\tau - \tau_l)
  \tag{2-2}
\end{equation}
式中：$\alpha_l$ 与 $\varphi_l$ 分别为第 $l$ 条多径的幅度与相位；$\tau_l$ 为时延；$L$ 为多径数量；$\delta(\cdot)$ 为狄拉克 $\delta$ 函数。该形式将信道表征为若干时延抽头的叠加，便于理解多径结构与人体活动对路径时延、衰减的调制作用。

按输出类型，基于 CSI 的 Wi-Fi 感知应用可分为三类：检测（二分类，如人体存在检测、跌倒检测）、识别（多分类，如活动识别、手势识别、行人识别）与估计（数值输出，如呼吸率估计、定位）\cite{ma2019wifi}。本文关注的室内动作识别属于识别类应用。Wi-Fi 感知在实际部署中面临共性挑战：鲁棒性与泛化（不同环境、设备、人体差异）、隐私与安全（无线信号可被窃听）、传感与通信共存（感知与数据传输共享信道资源）\cite{ma2019wifi}。

综上，基于 Wi-Fi CSI 的人类活动识别原理可概括为：利用 CSI 对多径信道幅度与相位的精细刻画，建立人体行为与环境信道变化之间的映射；通过信号预处理、特征提取与分类或估计算法，从 CSI 数据中推断室内动作类别。该原理为本文第三、四章中面向 Wi-Fi CSI 的类增量学习与少样本类增量学习方法提供了数据基础与物理背景。在本文后续实验中，原始 CSI 经预处理后组织为固定长度的时序特征序列（如 270 维 $\times$ 1000 时间步），作为第三、四章模型的输入。因而，本文后续方法设计重点不再放在信号级去噪，而放在基于时序表示的分类模型构建。


\xsection{增量学习}{Incremental Learning}

类增量学习关注的是类别集合随时间不断扩展时模型应如何完成持续更新。与一次性获得全部类别数据的离线训练不同，该任务要求模型在后续阶段仅利用当前到达的数据学习新类别，同时继续保持对历史类别的识别能力。其困难不只在于新旧知识之间存在竞争关系，还在于模型容量、样本分布和训练目标会随着增量过程不断变化。按照测试条件和增量对象的不同，持续学习通常可区分为类增量学习、任务增量学习和域增量学习等设定；其中，类增量学习在推理阶段不提供任务标识，需要模型直接在全部已见类别范围内完成判别，因此约束更强，也更适合刻画本文所关注的类别持续扩展场景\cite{zhou2024cil}。

除数据重放、动态网络与知识蒸馏外，还有基于数据正则化（如 GEM\cite{lopezpaz2017gem} 约束梯度方向）与参数正则化（如 EWC\cite{kirkpatrick2017ewc} 约束重要参数漂移）的方法；但在 CIL 场景下，这两类方法的效果多不如回放与动态网络\cite{zhou2024cil}。鉴于第三章方法主要由记忆回放、分类器扩展和蒸馏巩固三部分组成，本节重点介绍与上述设计直接相关的数据重放、动态网络和知识蒸馏方法。

\xsubsection{数据重放}{Data Replay}

在类增量学习中，数据重放的基本思路是为模型保留一部分历史信息，使其在适应新类别时不至于完全脱离旧类别分布。最常见的做法是在每次增量训练时，将当前任务样本与旧类别中选取的代表样本共同用于参数更新。这样一来，模型在优化新类别判别边界的同时，仍会持续接触旧类别特征，从而减轻旧知识被新知识覆盖的现象。由于这一策略直接作用于训练数据层面，其实现方式相对直观，也因此成为类增量学习中使用最广泛的基础手段之一\cite{zhou2024cil}。

该思路与人类认知中的“复习”现象一致：在接触新知识时反复回顾旧知识有助于巩固记忆\cite{zhou2024cil}。

直接回放（Direct Replay）是最直观的形式。模型维护一个 exemplar 集合 $\mathcal{E}$，在每轮增量更新时将当前任务数据 $\mathcal{D}^b$ 与 exemplar 联合训练，其损失可写为
\begin{equation}
  \mathcal{L} = \sum_{(\mathbf{x}, y) \in (\mathcal{D}^b \cup \mathcal{E})} \ell(f(\mathbf{x}), y),
  \tag{2-3}
\end{equation}
其中 $\ell(\cdot, \cdot)$ 为交叉熵等损失函数。相较仅使用 $\mathcal{D}^b$ 的基线 Finetune，该形式使模型在学习新类时持续回顾旧类样本，从而减弱遗忘。

Exemplar 的选取与管理直接影响重放效果。Herding 策略迭代选取最接近类均值的样本，使选中的 exemplar 平均特征尽量贴近类中心；此外还有基于预测熵、决策边界距离、数据增强不确定性等策略。在内存约束下，常采用固定总容量 $M$、每类分配 $\lfloor M / |\mathcal{V}_b| \rfloor$ 个样本的方式，避免随任务数线性增长\cite{zhou2024cil}。为降低存储成本，部分工作采用特征回放（存储特征向量而非原始图像）或低分辨率图像回放，但需额外适配以克服分布偏移。

生成式回放（Generative Replay）利用生成模型（如 GAN、VAE、扩散模型）合成旧类伪样本，从而无需存储原始 exemplar。该类方法依赖生成质量，在简单数据集上效果尚可，但在复杂大规模数据上常因生成失真而性能受限；且生成模型本身也会发生灾难性遗忘，需在分类器与生成器两端同时设计抗遗忘机制\cite{zhou2024cil}。

直接回放在 CIL 中应用广泛，但面临过拟合、新旧类数据不平衡以及隐私等问题。针对过拟合，可通过对 exemplar 施加 Lipschitz 约束或平衡采样加以缓解；在隐私或存储受限场景下，则需采用无 exemplar 或隐私友好的设计。本文第三章的类增量学习方法即采用数据重放策略，通过类均值最近邻迭代选择构建记忆库，缓解旧类数据不可及问题。

\xsubsection{动态网络}{Dynamic Networks}

与依靠有限样本回顾旧知识的思路不同，动态网络更强调通过扩展模型结构来为新类别分配额外表示能力。其核心出发点在于：当固定容量网络反复适配新任务时，新学习到的特征往往会挤占原有表征空间，从而引发遗忘；若能够在增量过程中为新任务引入新的参数模块，则旧知识可以在原有结构中相对稳定地保留下来，而新知识则由新增结构承担。这类方法并不试图完全依赖单一参数集合兼顾全部阶段，而是通过结构扩展在稳定性与可塑性之间取得平衡\cite{zhou2024cil}。

神经元扩展早期工作（如 DEN\cite{yoon2018den}）采用选择、扩展、复制与消除等流程：面对新任务时先选择性重训相关神经元，若损失仍高则自上而下扩展新神经元，并通过组稀疏正则剔除冗余。RCL、NAS 等将扩展过程形式化为强化学习或结构搜索问题。骨干扩展则进一步复制整条骨干网络。PNN\cite{rusu2016pnn} 为每个任务学习新骨干并冻结旧骨干，通过层间连接复用旧特征。DER\cite{yan2021der} 在每任务扩展新骨干，将新旧骨干特征拼接后接入更大的全连接层，输出可写为 $f(\mathbf{x}) = W_{\mathrm{new}}^\top [\phi_{\mathrm{old}}(\mathbf{x}), \phi_{\mathrm{new}}(\mathbf{x})]$，旧骨干冻结以保持旧类表征。该设计带来参数量线性增长，FOSTER 在扩展后引入知识蒸馏压缩阶段，将双骨干蒸馏回单骨干，在控制参数量的同时保留判别能力\cite{wang2022foster}。MEMO\cite{zhou2023memo} 将骨干解耦为共享浅层与任务专属深层，仅扩展深层块，在内存约束下实现高效扩展。

在 ViT 场景下，Prompt 扩展（如 L2P\cite{wang2022l2p}、DualPrompt\cite{wang2022dualprompt}、DyTox\cite{douillard2022dytox}）通过为每任务扩展 task token 或 prompt，以较轻量方式适配新类，多依赖预训练骨干。本文第三章的方法采用“分类器扩展 + 知识蒸馏巩固”的思路，与 FOSTER 的扩充—压缩范式在设计理念上相呼应。

\xsubsection{知识蒸馏}{Knowledge Distillation}

知识蒸馏在类增量学习中的作用，主要是利用上一阶段模型所保留的判别信息约束当前模型的更新方向。与仅依赖标签监督不同，蒸馏方法会进一步要求新模型在旧类别相关输出、特征表示或样本关系上尽量接近旧模型，从而在学习新类别时减少旧知识的大幅偏移。由于蒸馏不一定依赖大量历史原始数据，因此在存储受限或回放样本较少的条件下具有较高实用性，也是许多增量学习方法中常见的组成部分\cite{zhou2024cil,hinton2015distilling}。

将上一阶段模型 $f^{b-1}$ 视为教师、当前更新模型 $f$ 视为学生时，可通过约束二者在旧类上的输出或特征一致性将旧知识迁移至新模型\cite{hinton2015distilling}。

Logit 蒸馏约束旧类输出概率一致。LwF\cite{li2017lwf} 在交叉熵损失外添加蒸馏正则项，使新模型在旧类上的软标签与旧模型一致。iCaRL\cite{rebuffi2017icarl} 结合 exemplar 与蒸馏，并采用最近类均值推理。BiC\cite{wu2019bic} 引入动态权衡系数 $\lambda = |\mathcal{V}_{b-1}| / |\mathcal{V}_b|$，并附加偏差校正层以缓解新旧类 logit 偏移。Feature 蒸馏则约束嵌入或中间特征。UCIR\cite{hou2019ucir} 使用余弦相似度约束新旧嵌入一致；LwM\cite{dhar2019lwm} 约束注意力图；PODNet\cite{douillard2020podnet} 对池化后的空间特征做蒸馏。Relational 蒸馏关注样本间结构关系：R-DFCIL 对齐三元组角度关系；TPCIL、TOPIC 等通过图或神经气网络建模类间关系并对齐。

蒸馏常配合温度软化 softmax 与 KL 散度使用。其难点在于稳定性—可塑性权衡：蒸馏权重过大将抑制新类学习，过小则难以抵抗遗忘。与动态网络相比，蒸馏方法在内存受限时往往更优，但在充足内存下动态网络的扩展表征能力更具优势\cite{zhou2024cil}。

图~\ref{fig:cil-kd} 给出了知识蒸馏在 CIL 中的三类形式示意。综上，数据重放、动态网络与知识蒸馏可单独或组合使用，为本文后续面向 Wi-Fi CSI 的类增量学习设计提供了方法基础。

\begin{figure}[H]
  \centering
  \includegraphics[width=0.85\textwidth]{fig2-2-cil-kd.png}
  \caption{CIL 中知识蒸馏的三种形式}
  \label{fig:cil-kd}
\end{figure}


\xsection{少样本学习}{Few-Shot Learning}

少样本学习关注的是在监督信息有限的条件下如何获得可迁移的判别能力。与依赖大量标注样本拟合类别边界的常规监督学习不同，少样本学习更强调通过任务结构、相似度度量或先验知识提高模型对少量样本的利用效率\cite{tian2024fscilsurvey,zhaokailin2021fslreview}。在典型的 $N$-way $K$-shot 设定中，模型只能借助每类极少量支持样本完成对查询样本的分类，因此其性能很大程度上取决于表示空间是否具有足够的泛化性，以及模型是否能够从有限样本中提取稳定类别信息\cite{tian2024fscilsurvey,zhaokailin2021fslreview}。考虑到第四章方法需要在极少样本条件下完成新类建模，本节重点介绍与原型表示、相似度度量和样本扩展直接相关的少样本学习方法，为后续双分类器与图注意力建模提供背景。

\xsubsection{数据增强}{Data Augmentation}

数据增强通过增加已有数据的多样性缓解样本稀缺，降低过拟合风险并提升泛化能力\cite{tian2024fscilsurvey}。按数据来源可分为三类：基于训练集变换、基于弱标签或无标签数据、基于相似数据集。第一类对支持集样本施加几何变换（旋转、裁剪、翻转）、颜色抖动等；第二类利用无标签或半监督数据扩展表示学习；第三类从语义相近的数据集迁移样本或特征。

除直接变换外，还可利用生成模型合成新样本或特征。VAE、GAN 等通过建模数据分布，在隐空间采样生成伪样本，或直接在特征空间生成合成特征，以扩充支持集并缓解过拟合\cite{tian2024fscilsurvey}。在 FSCIL 场景下，数据增强常与度量学习或原型学习配合，用于基础阶段增强或增量阶段的虚拟样本生成。

\xsubsection{度量学习}{Metric Learning}

度量学习方法在嵌入空间中通过样本与支持集之间的相似度或距离进行分类，无需大量数据即可优化度量以拉近同类样本、拉远异类样本\cite{tian2024fscilsurvey}。

孪生网络（Siamese Network）\cite{koch2015siamese}采用双分支结构，对一对输入样本分别提取特征并计算特征间距离，通过对比损失学习判别性嵌入。匹配网络（Matching Network）\cite{vinyals2016matching}利用注意力机制对支持集与查询样本进行动态匹配与聚合，为每个查询生成与类别相关的表示。原型网络（Prototypical Network）将每类表征为其支持集样本嵌入的均值，即类原型 $\boldsymbol{\mu}_c = \frac{1}{|S_c|} \sum_{\mathbf{x} \in S_c} f_\theta(\mathbf{x})$；推理时根据查询嵌入与各原型的欧氏距离或余弦相似度进行最近邻分类。图~\ref{fig:fsl-metric} 给出了上述三种网络的结构对比。

度量学习的核心在于原型表示与相似度度量的设计。在 FSCIL 中，原型网络因结构简洁、易于扩展新类而被广泛采用；但其在增量场景下面临原型漂移、新旧类冲突等问题。本文第四章在原型表示基础上引入图注意力网络建模类间关系，并采用双分类器架构缓解类别级过拟合，其原型初始化与余弦分类的设计与度量学习及原型学习的思路一脉相承。

除度量学习外，FSL 还有模型类与优化类方法。模型类方法通过特定网络结构应对少样本：记忆增强网络（MANN）用外部记忆显式存储类别信息，利用长时记忆进行快速学习；Meta 网络学习跨任务的元级知识，通过快速参数化调整归纳偏置以实现泛化\cite{tian2024fscilsurvey}。优化类方法侧重训练策略：模型无关元学习（MAML）\cite{finn2017maml}在多任务上训练初始参数，使新任务上能快速收敛；Reptile\cite{nichol2018reptile} 将 MAML 的二步梯度简化为单步，降低计算开销。这些方法与度量学习互补，在 FSCIL 中常与原型、图结构等组合使用。

\begin{figure}[H]
  \centering
  \includegraphics[width=0.85\textwidth]{fig2-3-fsl-metric.png}
  \caption{度量学习方法常见网络结构}
  \label{fig:fsl-metric}
\end{figure}


\xsection{少样本类增量学习}{Few-Shot Class-Incremental Learning}

少样本类增量学习同时受到类别持续扩展与新类样本极少这两重约束\cite{tian2024fscilsurvey,tao2020topic}。模型不仅要在后续会话中识别此前未见的新类别，还要在旧类样本难以访问的条件下尽量维持已有类别的判别边界。相比第三章讨论的类增量学习，FSCIL 的困难进一步体现在：新类样本不足会削弱类别原型或决策边界的稳定性，而持续更新又会放大新旧类别之间的竞争关系。因此，该任务本质上要求模型同时处理表示学习不足、类别偏置和灾难性遗忘等多个问题\cite{tian2024fscilsurvey,tao2020topic}。在中文研究中，姚光乐等提出了基于特征分布学习的小样本类增量学习方法，通过利用已学习类别的特征分布信息来表示知识，并在增量阶段融合新类特征分布，从而缓解少样本条件下的新类过拟合与旧类遗忘问题\cite{yaoguangle2023fscildist}。

现有 FSCIL 方法通常采用「先在基础类别上学习通用表征，再在增量阶段适配新类别」的总体训练思路。其原因在于，基础阶段拥有相对充足的数据，更适合用于构建稳定的特征提取器；而进入增量阶段后，模型需要在极少样本条件下快速完成新类接入，因此更新重点往往会从完整表征学习转向分类器调整、原型构建或关系建模。围绕这一基本流程，不同研究又在数据利用方式、网络结构设计和训练目标构造上进行了进一步扩展\cite{tian2024fscilsurvey}。为缓解少样本下的过拟合与遗忘，部分工作扩展监督策略：在增量会话中引入无标签样本，将半监督学习与 FSCIL 结合；或利用自监督预训练获得双骨干特征提取器，再通过高斯生成器与特征融合适应新类\cite{tian2024fscilsurvey}。本节从基于数据、基于结构与基于优化三个维度概述 FSCIL 主流方法，为本文第四章少样本增量学习方法的设计提供理论背景。

\xsubsection{基于数据}{Data-Based Methods}

基于数据的 FSCIL 方法通过重放或合成旧类样本，使模型在学习新类时仍能接触旧类信息。直接重放维护 episode 记忆 $\mathcal{M}$，在增量更新时将当前任务数据 $D^{(t)}$ 与 $\mathcal{M}$ 联合训练，损失形式与式（2-3）类似。但直接存储 exemplar 在少样本场景下面临存储、样本选择与隐私约束；且新旧类样本量失衡时，知识蒸馏等依赖 exemplar 的方法易在性能权衡上偏于旧类\cite{tian2024fscilsurvey}。

生成式重放利用生成模型合成旧类伪样本，避免存储原始数据。数据无重放（data-free replay）方案通过熵正则鼓励生成器产生接近决策边界的“不确定”样本；FSILGAN 等针对少样本场景设计 GAN，将类别特定的合成视觉特征与语义向量对齐，并利用抗模式塌缩正则保证合成特征的多样性与可区分性。生成式重放的挑战在于少样本下难以逼近真实数据分布，且生成质量直接影响抗遗忘效果\cite{tian2024fscilsurvey}。

除真实或合成重放外，部分工作通过构建伪增量场景进行元学习。例如从基础类中采样模拟增量任务序列，在元训练阶段学习适应新知识与保留旧知识的元目标；LIMIT\cite{zhou2022limit} 等在基础集上构造“假”FSCIL 任务，通过不同假任务的元学习获得更具泛化性的特征，为真实增量阶段做准备\cite{tian2024fscilsurvey}。

\xsubsection{基于结构}{Structure-Based Methods}

基于结构的方法通过动态调整网络拓扑或模块，使模型随增量任务扩展表示能力。Tao 等提出的 TOPIC\cite{tao2020topic} 利用神经气（Neural Gas, NG）网络学习特征空间中各类形成的拓扑结构；通过保持 NG 拓扑稳定性抵抗遗忘，同时随新样本动态扩展以改善少样本新类表示。Zhang 等提出的 CEC\cite{zhang2021cec} 将训练解耦为嵌入学习与分类器学习：先在基础数据上训练骨干获得强特征提取能力，再在分类器层引入图注意力模型；随增量任务到达，图模型的节点与边权动态增长，使分类器适应全部类别。Zou 等发现大间隔分类虽有利于基础类，但易导致学习新类时性能下降，称为类别级过拟合；CLOM\cite{zou2022clom} 引入基于边界的额外约束以缓解该问题，并采用双分类器架构分离基础类与新类决策。

上述图注意力与双分类器思路与本文第四章少样本增量学习方法一脉相承：该方法在原型基础上引入增强图注意力建模类间与时序关系，并采用双分类器缓解类别级过拟合。

\xsubsection{基于优化}{Optimization-Based Methods}

基于优化的方法从表示学习、知识蒸馏与目标函数形态等方面调整训练策略。表示学习方面，前向兼容训练（FACT）\cite{zhou2022fact}在特征空间预分配虚拟原型，为未来新类预留空间，通过实例混合生成虚拟样本使模型具有前向兼容性；LIMIT\cite{zhou2022limit} 等则通过元校准模块生成实例特定的嵌入以提升泛化。知识蒸馏在 FSCIL 中可用于转移旧知识，但少样本场景下新旧类失衡使其效果受限；语义感知蒸馏通过引入词嵌入等辅助信息、图关系蒸馏（如 ERDIL）通过维护类间关系图进行知识抽取与表示，一定程度上缓解了上述问题\cite{tian2024fscilsurvey}。

从函数优化角度，F2M\cite{shi2021f2m} 等工作发现基础训练所得的平坦局部极小值比尖锐极小值具有更好的泛化能力；在平坦区域内对新任务进行微调可大幅减轻灾难性遗忘。稳定性—可塑性权衡贯穿 FSCIL：过度稳定将抑制新类学习，过度可塑则加速遗忘；ERM 在样本稀缺时不可靠，需结合元学习、正则化或结构设计加以应对。与视觉图像不同，Wi-Fi CSI 数据具有显著的时序性和跨子载波结构，因此现有 FSCIL 方法在迁移到 Wi-Fi 感知任务时需考虑时序表征与模态适配问题。综上，基于数据、结构与优化的方法可单独或组合使用，为面向 Wi-Fi CSI 的少样本类增量学习（如本文第四章）提供方法基础。


\xsection{本章小结}{Chapter Summary}

综上，本章从 Wi-Fi CSI 感知原理、类增量学习、少样本学习和少样本类增量学习四个方面梳理了本文所需的理论基础。其中，CSI 时序表征构成了第三、四章共同的数据基础；数据重放、动态扩展和知识蒸馏构成了第三章方法设计的核心依据；原型表示、相似度建模和类间关系传播则为第四章方法设计提供了直接支撑。基于此，下一章将首先讨论样本相对充足条件下的类增量学习方法。
